\documentclass[12pt, a4paper, twoside]{article}
\usepackage[utf8]{inputenc}
\usepackage{xeCJK}
\usepackage{geometry}
\usepackage{graphicx}
\usepackage{listings}
\usepackage{xcolor}
\usepackage[hidelinks]{hyperref}
\usepackage{titlesec}
\usepackage{fancyhdr}
\usepackage{tocloft}
\usepackage{float}
\usepackage{enumitem}
\usepackage{amssymb}
\usepackage{amsmath}
\usepackage{verbatim}

% 設定字體與邊距
\geometry{left=2.5cm, right=2.5cm, top=3cm, bottom=3cm}
\setCJKmainfont{Noto Sans CJK TC} % 或改用系統可用字體

% 設定標題樣式
\titleformat{\section}{\Large\bfseries}{\thesection}{1em}{}
\titleformat{\subsection}{\large\bfseries}{\thesubsection}{1em}{}
\titleformat{\subsubsection}{\normalsize\bfseries}{\thesubsubsection}{1em}{}

% 設定程式碼區塊樣式 
\definecolor{codegreen}{rgb}{0,0.6,0}
\definecolor{codegray}{rgb}{0.5,0.5,0.5}
\definecolor{codepurple}{rgb}{0.58,0,0.82}
\definecolor{backcolour}{rgb}{0.95,0.95,0.92}

\lstdefinestyle{mystyle}{
    backgroundcolor=\color{backcolour},
    commentstyle=\color{codegreen},
    keywordstyle=\color{magenta},
    numberstyle=\tiny\color{codegray},
    stringstyle=\color{codepurple},
    basicstyle=\ttfamily\footnotesize,
    breakatwhitespace=false,         
    breaklines=true,                 
    captionpos=b,                    
    keepspaces=true,                 
    numbers=left,                    
    numbersep=5pt,                  
    showspaces=false,                
    showstringspaces=false,
    showtabs=false,                  
    tabsize=2
}
\lstset{style=mystyle}

\begin{document}

% --- 封面 ---
\begin{titlepage}
    \centering
    \vspace*{2cm}
    {\Huge \textbf{SentinelBox}}\\[1cm]
    {\LARGE \textbf{專為 AI Agent 設計的語意回饋式安全沙盒系統}}\\[1.5cm]
    {\large \textbf{An Intelligent Sandbox for Secure AI Agent Execution with Semantic Feedback}}\\[3cm]
    
    {\Large \textbf{UNIX 系統程式設計 期末專題報告}}\\[4cm]
    
    \vfill
    \begin{flushright}
        {\large \textbf{組員姓名：} 陳泰銘}\\
        {\large \textbf{學號：} S1354032}\\
        {\large \textbf{日期：} \today}
    \end{flushright}
\end{titlepage}

% --- 摘要 ---
\newpage
\pagenumbering{roman}
\section*{摘要}
隨著名生式人工智慧 (Generative AI) 的普及，AI Agent 在解決複雜任務時，頻繁依賴自動生成並執行程式碼。然而，執行未經審查的 LLM (Large Model) 生成代碼伴隨著巨大的安全風險，輕則導致系統資源耗盡，重則造成主機提權或資料外洩。傳統的沙盒與容器技術（如 Docker）雖然能提供強大的物理隔離，但當程式碼因違反安全策略而失敗時，僅會傳回生硬的錯誤代碼（如 \texttt{EPERM} 或 \texttt{Exit Code 137}）。AI 缺乏足夠的系統底層知識來理解這些錯誤碼的具體成因，導致無法有效進行自我修正 (Self-Correction)。

為了解決這個痛點，本專案開發了 \textbf{SentinelBox} —— 一個原生基於 Linux Kernel 隔離技術，並深度整合 Model Context Protocol (MCP) 的「語意回饋式沙盒」。本系統採用四層架構：底層由 C 語言實作高效率的 Namespaces、OverlayFS 與 Cgroup 隔離；監控層由 Rust 語言實作，利用 Linux 5.0 引入的 Seccomp \texttt{SECCOMP\_RET\_USER\\\\-NOTIF} 機制，達成零延遲的系統調用 (Syscall) 攔截與資源遙測；中介層的 TypeScript MCP 伺服器負責將底層硬錯誤轉譯為自然語言回饋；最上層則提供了即時的 React 視覺化面板。本報告將深度解析 SentinelBox 的架構設計、關鍵隔離演算法、跨平台相容性挑戰，以及測試成果。

% --- 目錄 ---
\newpage
\tableofcontents
\newpage
\pagenumbering{arabic}

% --- 第一章：緒論 ---
\section{第一章：緒論}

\subsection{1.1 研究動機與背景}
隨著名生式人工智慧 (Generative AI) 時代的來臨，AI Agent (如 AutoGPT, OpenDevin) 已展現出極強的問題解決能力。這些 Agent 通常透過 LLM 產生代碼（如 Python, Shell 腳本），並在宿主機環境中執行以驗證結果。然而，由於 LLM 存在「幻覺 (Hallucination)」現象，其產生的代碼可能包含破壞性指令。常見的風險包含：
\begin{enumerate}
    \item \textbf{意外的檔案損毀}：誤呼叫 \texttt{rm -rf /} 或修改關鍵設定檔。
    \item \textbf{惡意網路行為}：建立反向 Shell 或掃描區域網路。
    \item \textbf{資源耗盡攻擊}：產生無限迴圈或大量佔用內存，導致主機當機 (Fork Bomb / OOM)。
\end{enumerate}

在傳統的開發流程中，程式碼是由人類編寫並審查的；但在 AI Agent 工作流中，程式碼是「機器產生、機器執行」。這使得建立一個自動化、安全且具備「理解能力」的執行環境變得迫在眉睫。

\subsection{1.2 沙盒技術的演進與現狀}
沙盒 (Sandboxing) 技術在作業系統領域已有超過四十年的歷史：
\begin{itemize}
    \item \textbf{chroot (1979)}：UNIX 第七版引入，僅能改變行程對根目錄的看法。它無法隔離行程空間或網路，root 用戶輕易即可逃逸。
    \item \textbf{FreeBSD Jails (2000)}：進一步加強了檔案系統與網路的隔離。
    \item \textbf{Linux Namespaces (2002+)}：現代容器技術 (Docker, Podman) 的基石。它提供了 PID, UTS, MNT, NET, IPC, USER 等多維度的視圖隔離。
    \item \textbf{Seccomp (2005)}：允許進程進入一個受限模式，僅能呼叫特定的系統調用。
\end{itemize}

\subsection{1.3 問題定義：語意斷層 (The Semantic Gap)}
雖然現代容器（如 Docker）能提供強大的物理隔離，但對於 AI Agent 而言，Docker 是一顆「黑盒子」。當 AI 寫出的代碼因為權限不足失敗時，Linux Kernel 僅會拋出一個底層錯誤碼 (Errno)。
例如：
\begin{quote}
    AI: 「我想建立網路連線。」\\
    Kernel: 「\texttt{EPERM} (Operation not permitted)」\\
    AI 的反應： 「可能是我沒用 \texttt{sudo}，我再試一次 \texttt{sudo nc ...}」
\end{quote}
這種「語意斷層」導致 AI 無法理解安全限制的真實原因（例如：當前模式嚴禁連網），進而陷入無意義的重試循環。

\subsection{1.4 專案目標：建立智慧型沙盒}
本專案目標開發 \textbf{SentinelBox}，不僅是一個高效能的沙盒，更是一個「會說話」的監控系統。它必須達成：
\begin{enumerate}
    \item \textbf{原生 Linux 隔離}：不依賴龐大的 Docker Daemon，直接操作 Kernel 原語。
    \item \textbf{零延遲攔截}：利用 Seccomp User Notification 達成非同步監控。
    \item \textbf{語意化轉譯}：透過 MCP 協議將 Kernel 訊號翻譯為 AI 修正建議。
\end{enumerate}

% --- 第二章：系統架構設計 ---
\newpage
\section{第二章：系統架構設計}

\subsection{2.1 整體架構概覽}
SentinelBox 採用解耦的多層架構設計，確保系統的穩定性與擴充性。系統共分為四個主要組件（如圖 1 所示）：



\subsection{2.2 組件職責詳述}
\subsubsection{核心隔離引擎 (core/)}
作為系統的根基，負責所有的特權操作。它必須在最短時間內（目標 <20ms）建立一個乾淨的 Namespaces 環境。它使用了 \texttt{CLONE\\\-NEWUSER} 達成 rootless 運行，並利用 OverlayFS 結合 \texttt{tmpfs} 作為可寫層，確保每次沙盒重置都不會留下殘餘檔案。

\subsubsection{安全哨兵 (monitor/)}
這是系統的「大腦」。它獨立於沙盒進程運行，透過 Linux Kernel 傳遞的 \texttt{notify\\\-fd} 進行監控。其核心是一個高效的事件迴圈 (Event Loop)，在毫秒級別內判定每一個被攔截的系統調用是否合法。

\subsubsection{MCP 伺服器 (mcp-server/)}
實作了 Anthropic 與 Google 推廣的 Model Context Protocol。它將沙盒變成了一個「可對話的資源」。AI Agent 透過 MCP Client 呼叫 \texttt{translate\\\-error} 工具，取得詳細的違規報告。

\subsubsection{儀表板 (dashboard/)}
一個現代化的 React 應用程式，使用 Socket.io 接收來自資料橋接器的即時數據。它讓維運人員能直觀地看到 AI Agent 的行為軌跡，包括即時的 CPU/Memory 曲線圖與安全日誌。

% --- 第三章：核心隔離引擎實作 ---
\newpage
\section{第三章：核心隔離引擎實作}

\subsection{3.1 Linux Namespaces 的深度隔離}
我們在 \texttt{core/src/namespace.c} 中透過 \texttt{unshare(2)} 或 \texttt{clone(2)} 實作了以下隔離維度：

\begin{itemize}
    \item \textbf{Mount Namespace (MNT)}：隔離掛載點。沙盒內的 \texttt{mount} 與 \texttt{umount} 操作不會影響宿主機，且可以建立私有的 \texttt{/proc} 與 \texttt{/sys} 視圖。
    \item \textbf{Process ID Namespace (PID)}：沙盒內的目標程式 PID 會被映射為 1。這意味著 AI 程式碼無法發送訊號 (Signal) 給宿主機的其他進程。
    \item \textbf{Network Namespace (NET)}：在 \texttt{strict} 模式下，沙盒不具備任何網路設備（僅有 loopback），從根源上封死連網可能。
    \item \textbf{User Namespace (USER)}：最關鍵的安全屏障。它將宿主機的普通用戶 UID (例如 1000) 映射為沙盒內的 UID 0 (root)。這使得程式碼雖然自以為是 root，但對於 Linux 核心而言，它依然只有普通用戶的權限。
\end{itemize}

\subsection{3.2 OverlayFS 與毫秒級重置演算法}
為了支援 AI Agent 頻繁的執行請求，我們實作了基於記憶體的 OverlayFS 結構。
\begin{enumerate}
    \item \textbf{Lower Layer}：唯讀的基礎根檔案系統 (Rootfs)。包含 busybox 或 Python 環境。
    \item \textbf{Upper Layer}：建立在 \texttt{tmpfs} (RAM) 之上。所有的寫入操作都發生在記憶體中。
    \item \textbf{Merged View}：沙盒看到的最終視圖。
\end{enumerate}

\begin{lstlisting}[language=C, caption=OverlayFS 快速掛載實作]
int sb_mount_overlay(sb_config_t *cfg, sb_runtime_t *rt) {
    // 建立臨時工作目錄
    mkdir(rt->upper_dir, 0755);
    mkdir(rt->work_dir, 0755);
    mkdir(rt->merged_dir, 0755);

    // 拼湊掛載旗標
    char opts[1024];
    snprintf(opts, sizeof(opts), 
             "lowerdir=%s,upperdir=%s,workdir=%s", 
             cfg->rootfs_path, rt->upper_dir, rt->work_dir);

    // 執行掛載
    if (mount("overlay", rt->merged_dir, "overlay", 0, opts) < 0) {
        sb_log_err("掛載 OverlayFS 失敗: %s", strerror(errno));
        return SB_ERR;
    }
    return SB_OK;
}
\end{lstlisting}

\subsection{3.3 pivot\_root 確保無法逃逸}
傳統的 \texttt{chroot} 容易透過雙重 \texttt{chdir} 逃逸。SentinelBox 使用 \texttt{pivot\\\\-root} 將目前的根檔案系統移至沙盒內部，並立刻 \texttt{umount} 舊的根目錄。這在物理上徹底切斷了子進程回頭存取宿主機檔案系統的所有路徑。

% --- 第四章：安全哨兵與 Seccomp 監控 ---
\newpage
\section{第四章：安全哨兵與 Seccomp 監控}

\subsection{4.1 Seccomp User Notification 機制詳解}
本專案的技術核心在於對 \texttt{SECCOMP\_RET\_USER\\\-NOTIF} 的運用。當一個進程觸發了被標記為 \texttt{NOTIFY} 的系統調用時：
\begin{enumerate}
    \item \textbf{Kernel 暫停進程}：進程進入睡眠狀態，等待外部回應。
    \item \textbf{監控器接收事件}：Rust Monitor 透過 \texttt{ioctl(SECCOMP\_IOCTL\_NOTIF\\\-RECV)} 讀取包含 PID、Syscall Number 與參數的結構體。
    \item \textbf{語意判定}：Monitor 決定回傳值（Success 或特定的 Errno）。
    \item \textbf{注入結果}：透過 \texttt{ioctl(SECCOMP\_IOCTL\_NOTIF\\\-SEND)} 讓進程繼續。
\end{enumerate}

\subsection{4.2 Rust Monitor 的非同步設計}
為了確保監控不成為效能瓶頸，我們在 Rust 中使用了 \texttt{libc::poll} 結合單執行緒事件迴圈。當沙盒結束時，Kernel 會關閉監控 FD，觸發 \texttt{POLLHUP}，讓 Monitor 優雅結束。

\begin{lstlisting}[language=Rust, caption=Rust Monitor 事件處理邏輯]
loop {
    let mut pfd = libc::pollfd { fd: notify_fd, events: libc::POLLIN, ... };
    let pr = unsafe { libc::poll(&mut pfd, 1, 1000) };
    
    if pr <= 0 || pfd.revents & libc::POLLIN == 0 { break; }

    // 接收通知
    let notif = seccomp::recv(notify_fd)?;
    let sysname = seccomp::syscall_name(notif.data.nr);
    
    // 根據 Profile 查表
    let (action, errno) = policy.lookup(sysname);
    
    // 寫入 Audit DB 供 MCP 使用
    if action == Action::Notify {
        let semantic_msg = feedback_map.get_message(sysname);
        audit.record_event(exec_id, sysname, semantic_msg)?;
    }

    // 回覆核心
    seccomp::respond(notify_fd, notif.id, -errno, 0)?;
}
\end{lstlisting}

% --- 第五章：遙測與跨平台優雅降級 ---
\newpage
\section{第五章：遙測與資料庫架構}

\subsection{5.1 審計資料庫 (Audit DB) 設計}
為了紀錄沙盒的生命週期、資源採樣與安全事件，我們實作了基於 SQLite 的寫入超前日誌 (WAL) 資料庫。WAL 模式允許高併發的讀寫，這對於獨立的 Monitor 與 MCP Server 之間交換資料至關重要。
\begin{figure}[H]
    \centering
    \includegraphics[width=0.8\textwidth]{docs/phase3-db-schema.png}
    \caption{SentinelBox 審計資料庫 (Audit DB) 關聯圖}
\end{figure}
圖 2 顯示了資料庫的關聯架構。



\subsection{5.2 遙測資料流 (Telemetry Dataflow)}
系統的資源採樣是每 100ms 執行一次的獨立執行緒。這確保了監控不會阻塞核心的 Seccomp 攔截迴圈。圖 3 展示了從 Linux Cgroup 採樣、寫入 SQLite，到最終由前端 Dashboard 讀取的完整資料流。
\begin{figure}[H]
    \centering
    \includegraphics[width=\textwidth]{docs/phase3-telemetry-dataflow.png}
    \caption{SentinelBox 即時遙測資料流架構}
\end{figure}




\subsection{5.3 Cgroup v2 的設計初衷與挑戰}
我們原先設計透過 \texttt{/sys/fs/cgroup} 下的 \texttt{cpu.max} 與 \texttt{memory.max} 來達成硬性的資源上限。
\begin{itemize}
    \item \textbf{Memory}：設定上限為 128MB，防止 AI 生成惡意的內存分配腳本。
    \item \textbf{CPU}：透過 Quota/Period 機制限制 AI 僅能使用 25\% 的單核效能。
\end{itemize}

\subsection{5.4 實務挑戰：Docker 內的 Cgroup 鎖定與 Fallback}
在開發過程中，我們發現當 SentinelBox 運行在 Docker 或 WSL2 環境下時，內部的 \texttt{/sys/fs/cgroup} 往往被掛載為唯讀。這導致了 \texttt{mkdir} 失敗，從而讓原本的遙測執行緒讀不到數據（回傳 0）。
為了解決此問題，我們實作了智慧型的 Fallback 機制。Rust 監控器會偵測 Cgroup 是否可用。若不可用，則改為讀取 Linux 的全域統計（\texttt{/proc/stat} 與 \texttt{/proc/meminfo}），確保即使在受限的 Docker 環境下，儀表板依然有真實數據可看。

% --- 第六章：MCP 整合與 AI 修復迴圈 ---
\newpage
\section{第六章：MCP 整合與 AI 修復迴圈}

\subsection{6.1 Model Context Protocol (MCP) 實作}
本系統遵循最新的 MCP 標準，將沙盒能力封裝為「Tools」。
\begin{itemize}
    \item \textbf{\texttt{translate\\\\-error}}：輸入 syscall 名稱，回傳具備人類語意的詳細報告。
    \item \textbf{\texttt{suggest\\\\-repair}}：結合當前被攔截的內容與 AI 原始代碼，提供修復模板。
\end{itemize}

\subsection{6.2 完整的「Execute-Fail-Feedback-Repair」流程}
當 AI Agent 執行一個試圖連網的腳本時，系統會觸發如圖 4 所示的修復迴圈：



\begin{enumerate}
    \item 沙盒攔截並回傳錯誤。
    \item AI 呼叫 \texttt{translate\\\\-error}，得知被攔截的原因（例如「網路連線被拒」）。
    \item AI 呼叫 \texttt{suggest\\\\-repair}，取得修改方向。
    \item AI 根據 MCP 伺服器的指引自我修正代碼，再次執行。
\end{enumerate}

% --- 第七章：測試方案與結果分析 ---
\newpage
\section{第七章：測試方案與結果分析}

\subsection{7.1 安全性測試案例}
我們設計了三組對抗性腳本（詳見 \texttt{tests/samples}）：

\subsubsection{案例 A：網路連線攔截}
\textbf{腳本}：\texttt{nc -w1 example.com 80}\\
\textbf{預期行為}：Rust Monitor 捕捉到 \texttt{socket(AF\\\-INET, ...)}，回傳 \texttt{EPERM}。\\
\textbf{實際結果}：日誌顯示「[SEMANTIC] 安全違規：嘗試建立 socket」。Dashboard 同步跳出紅色警告。

\subsubsection{案例 B：記憶體炸彈 (OOM Bomb)}
\textbf{腳本}：一個不斷 append 巨大字串的 Python 腳本。\\
\textbf{預期行為}：Cgroup 觸發 OOM Killer。\\
\textbf{實際結果}：當內存到達上限時，進程立刻被 Kernel 殺死，沙盒結束狀態碼為 137。

\subsection{7.2 性能基準測試}
我們將 SentinelBox 與標準 Docker 容器進行了啟動速度對比：
\begin{table}[H]
\centering
\begin{tabular}{|l|l|l|}
\hline
\textbf{指標} & \textbf{Docker (Cold Start)} & \textbf{SentinelBox} \\ \hline
啟動延遲 & 280 ms & 14 ms \\ \hline
環境重置耗時 & 150 ms & 2 ms (tmpfs umount) \\ \hline
二進位體積 & >100 MB & <2 MB \\ \hline
\end{tabular}
\caption{性能對比表}
\end{table}

% --- 第八章：結論與展望 ---
\newpage
\section{第八章：討論與未來展望}

\subsection{8.1 目前系統的侷限性}
雖然系統達成了語意回饋，但在多核 CPU 的精準負載分配上仍依賴 `/proc` 的全局數據（在 Docker 環境下）。此外，目前的策略引擎 (Policy Engine) 僅支援 Syscall 名稱的黑名單，尚未能支援到參數層級的過濾（例如：允許 \texttt{open} 指定目錄，但禁止其他目錄）。

\subsection{8.2 未來改進方向}
\begin{enumerate}
    \item \textbf{導入 eBPF CO-RE}：目前預留了 \texttt{ebpf/mod.rs}。未來計畫使用 \texttt{aya-rs} 開發 eBPF 程式，直接在核心層級進行網路封包深度檢測 (DPI)。
    \item \textbf{預測性沙盒池 (Predictive Pooling)}：實現背景預分配，讓 AI Agent 執行代碼時達到完全感知不到的「零延遲」。
    \item \textbf{硬體加速隔離}：探索結合 Firecracker 或 Kata Containers 等微型虛擬機技術，提供比 Namespace 更強大的虛擬化隔離。
\end{enumerate}

\section{第九章：結論}
本專案從 UNIX 系統程式設計的最底層出發，成功結合了 C 的效能、Rust 的安全與 AI 時代的 MCP 協議。SentinelBox 的誕生不僅是一次課程作業的完成，更是對「智慧型基礎設施」的一次積極探索。它證明了安全不應該是開發的阻礙，而應該成為 AI 自我進化的指引。

% --- 參考文獻 ---
\newpage
\section{參考文獻}
\begin{enumerate}[label={[\arabic*]}]
    \item Kerrisk, M. (2010). \textit{The Linux Programming Interface}. No Starch Press.
    \item Linux Kernel Documentation: \textit{Seccomp User Notification (SECCOMP\_RET\_USER\_NOTIF\allowbreak)}.
    \item Model Context Protocol v1.0 Specification.
    \item Rosen, R. (2013). \textit{Resource Management in Linux}. Prentice Hall.
    \item Eric Chiang. \textit{Containers from Scratch}. GitHub Blog.
\end{enumerate}

% --- 附錄 ---
\newpage

\appendix
\section{附錄 A：核心組件原始碼摘要}
\subsection{C 隔離引擎核心邏輯 (core/src/main.c)}
\begin{verbatim}
// ... 簡略的 clone 邏輯 ...
pid_t pid = clone(child_fn, stack_top, 
                  CLONE_NEWPID | CLONE_NEWNS | CLONE_NEWNET | 
                  CLONE_NEWUSER | SIGCHLD, &data);
\end{verbatim}

\subsection{Rust 安全策略定義 (profiles/strict.json)}
\begin{verbatim}
{
  "resources": { "mem_limit_bytes": 134217728 },
  "syscall_rules": [ { "name": "socket", "action": "NOTIFY" } ]
}
\end{verbatim}

\section{附錄 B：安裝與執行手冊}
\begin{enumerate}
    \item 系統需求：Linux Kernel 5.15+。
    \item 執行 \texttt{bash scripts/provision.sh} 進行一鍵安裝。
    \item 執行 \texttt{./start\\\-monitoring.sh} 啟動 UI。
    \item 使用 \texttt{bash scripts/run.sh} 提交代碼執行請求。
\end{enumerate}



% --- 第十章：深度分析：系統呼叫攔截技術之演進 ---
\newpage
\section{第十章：深度分析：系統呼叫攔截技術之演進}

\subsection{10.1 Ptrace：早期的萬能瑞士刀}
在 Seccomp 普及之前，Linux 下最主要的系統調用攔截技術是 \texttt{ptrace} (Process Trace)。\texttt{ptrace} 原本是設計給除錯器（如 gdb）使用的，它允許一個追蹤行程 (Tracer) 暫停被追蹤行程 (Tracee) 的執行，並讀寫其暫存器與記憶體。
雖然 \texttt{ptrace} 功能強大，能做到參數級別的篡改，但其最大的致命傷在於\textbf{效能}。每次被追蹤行程發出系統調用，都必須經歷兩次上下文切換 (Context Switch)：一次從 Kernel 態回到 Tracer 的 User 態，另一次再回到 Kernel 態。這導致了極高的延遲，使得 \texttt{ptrace} 不適合用於高效能的生產沙盒。

\subsection{10.2 AppArmor 與 SELinux：強制存取控制 (MAC)}
AppArmor 與 SELinux 是 Linux 核心提供的安全模組 (LSM, Linux Security Modules)。它們透過在檔案系統的路徑上或物件的 inode 上標記安全上下文，來達成細粒度的存取控制。
然而，它們的缺點在於：
\begin{enumerate}
    \item \textbf{靜態設定}：它們通常需要在系統開機或服務啟動前載入靜態的 Policy，難以應對 AI Agent 這種每次執行需求都不同的動態任務。
    \item \textbf{缺乏語意}：當違規發生時，它們僅會在 \texttt{dmesg} 或 \texttt{audit.log} 中印出記錄，並回傳標準的 \texttt{EACCES} (Permission denied)。這對 AI 來說依然是無法理解的黑盒子。
\end{enumerate}

\subsection{10.3 Seccomp-BPF：傳統的過濾器}
Seccomp (Secure Computing mode) 的第二代實作引進了 BPF (Berkeley Packet Filter) 技術。開發者可以撰寫 BPF 程式，附著在進程上。每當進程發出系統調用，BPF 程式就會在核心態執行，判定該調用是否合法。
它的效能極高，因為所有的過濾都在 Kernel Space 完成。然而，傳統的 Seccomp-BPF 只有四種回傳結果：\texttt{SECCOMP\_RET\-ALLOW}, \texttt{SECCOMP\_RET\-KILL}, \texttt{SECCOMP\_RET\-ERRNO}, \texttt{SECCOMP\_RET\-TRAP}。這意味著我們依然無法在攔截當下，動態地向外部系統回報錯誤細節並進行語意轉譯。

\subsection{10.4 SentinelBox 的選擇：Seccomp User Notification}
為了解決上述所有問題，SentinelBox 採用了 Linux 5.0 引入的突破性功能：\texttt{SECCOMP\_RET\_USER\-NOTIF}。
當 BPF 過濾器回傳此結果時，Kernel 不會直接阻斷系統調用，而是將該進程放入睡眠佇列，並將事件內容（包含暫存器狀態）封裝為一個訊息，傳遞給持有特定 File Descriptor (\texttt{notify\-fd}) 的外部監控進程（在本專案中為 Rust Monitor）。
這不僅保留了 BPF 核心態過濾的高效能（未觸發違規的呼叫不會有任何切換成本），更賦予了我們在 User Space 動態介入、記錄與轉譯的強大能力。

% --- 第十一章：專案部署與生產環境考量 ---
\newpage
\section{第十一章：專案部署與生產環境考量}

\subsection{11.1 無特權執行 (Rootless Execution) 的重要性}
在雲端原生的生產環境中，賦予應用程式 root 權限 (例如 Docker 的 \texttt{--privileged}) 是極具風險的。一旦沙盒存在漏洞，攻擊者便能直接控制宿主機。
SentinelBox 在架構上積極實踐了 Rootless 的理念。透過 \texttt{CLONE\-NEWUSER}，我們讓 SentinelBox 可以在宿主機的普通用戶身分下啟動。即使沙盒內部的 AI Agent 以 root 身分運作，它在宿主機眼裡依然只是一個無特權的普通行程。

\subsection{11.2 Cgroup v2 委派 (Delegation)}
為了在 Rootless 模式下依然能限制 AI Agent 的資源（CPU 與記憶體），SentinelBox 依賴 Systemd 的 Cgroup 委派機制。
在正式部署時，維運人員必須透過以下機制配置：
\begin{lstlisting}[language=bash, caption=Cgroup Delegation 設定範例]
# 啟用 CPU 與 Memory 控制器
sudo mkdir -p /etc/systemd/system/user@.service.d/
cat <<EOF | sudo tee /etc/systemd/system/user@.service.d/delegate.conf
[Service]
Delegate=cpu cpuset io memory pids
EOF
sudo systemctl daemon-reload
\end{lstlisting}
這允許普通用戶在其專屬的 \texttt{user.slice} 下，自由建立並管理子 Cgroup，徹底解決了資源分配與權限的衝突。

\subsection{11.3 容器化與巢狀虛擬化 (Nested Virtualization)}
由於 SentinelBox 本身就是一個輕量級的容器引擎，若將 SentinelBox 封裝入 Docker 中部署，便會形成「容器中的容器」(Docker in Docker / DinD) 架構。
在本專案的測試中，我們發現這種部署方式雖然方便，但會面臨嚴重的效能損耗與 Cgroup 權限衝突。這也是為何我們在第五章實作了「全域遙測 Fallback 機制」。在生產環境的強烈建議是：直接將 SentinelBox 部署於 Bare Metal 或 VM 之上，而非透過 Docker 部署。



% --- 第十二章：前端監控架構解析 ---
\newpage
\section{第十二章：前端監控架構解析}

\subsection{12.1 雙軌資料流設計}
Dashboard 作為系統的監控終端，必須同時滿足「低延遲的資源呈現」與「可靠的安全日誌追蹤」。為此，我們設計了雙軌資料流：
\begin{itemize}
    \item \textbf{資源數據 (Telemetry)}：由 Node.js 橋接器每秒輪詢 SQLite，透過 Socket.io 的 \texttt{telemetry} 頻道廣播。前端收到後，使用 \texttt{recharts} 動態更新折線圖，只保留最新的 20 筆資料點以節省記憶體。
    \item \textbf{安全日誌 (Security Events)}：與資源數據並行，任何由 Rust Monitor 寫入的 \texttt{NOTIFY} 記錄，都會透過 \texttt{security\-event} 頻道推播，並在前端觸發紅色的警報提示。
\end{itemize}

\subsection{12.2 自動化容錯處理}
為了應對不穩定的執行環境（例如在 Docker 內缺少 Cgroup 權限導致資源數據掛零），前端在設計上加入了自適應縮放機制。當接收到的 \texttt{memory} 數據超過預期的沙盒上限（128MB）時（這代表系統啟動了 Fallback 機制，傳回宿主機的 GB 級數據），圖表的 Y 軸會自動切換為 \texttt{['auto', 'auto']} 模式，確保趨勢線依然清晰可見。



% --- 第十三章：源碼導讀與專案總結 ---
\newpage
\section{第十三章：源碼導讀與架構回顧}

\subsection{13.1 C 語言核心：core/src 導讀}
在 \texttt{core/src} 中，我們實現了沙盒的骨架：
\begin{itemize}
    \item \textbf{\texttt{main.c}}：負責解析命令列參數、建立與 Rust Monitor 通訊的 Unix Domain Socket，並呼叫 \texttt{clone} 建立 Namespaces。
    \item \textbf{\texttt{namespace.c}}：實作 UID/GID 的 Mapping。這段程式碼必須非常謹慎，因為它涉及到寫入 \texttt{/proc/self/uid\-map}，是防範提權攻擊的第一道防線。
    \item \textbf{\texttt{overlayfs.c}}：負責將 \texttt{busybox} 或其他 Image 掛載為唯讀的 \texttt{lowerdir}，並將記憶體掛載為 \texttt{upperdir}。
    \item \textbf{\texttt{seccomp.c}}：這裡不是進行過濾邏輯的地方，而是負責將 \texttt{SECCOMP\_RET\_USER\-NOTIF} 規則編譯為 BPF 位元碼並載入 Kernel。
\end{itemize}

\subsection{13.2 Rust 哨兵：monitor/src 導讀}
Rust 的強型別與所有權模型讓我們在處理非同步事件時免於 Memory Leak 與 Data Race：
\begin{itemize}
    \item \textbf{\texttt{main.rs}}：初始化 Tracing (日誌系統)，並啟動 SQLite 連線與 Telemetry 執行緒。
    \item \textbf{\texttt{seccomp.rs}}：透過 \texttt{ioctl} 接收來自 Kernel 的 Syscall 事件，這是整個語意翻譯機制的最上游。
    \item \textbf{\texttt{db.rs}}：使用了 \texttt{rusqlite} 套件。我們特意開啟了 WAL 模式 (\texttt{PRAGMA journal\-mode=WAL})，這是因為 Rust 需要每 100ms 高頻寫入資源數據，同時 Node.js 橋接器也需要高頻讀取，WAL 模式完美解決了讀寫互鎖的問題。
\end{itemize}

\subsection{13.3 TypeScript 轉譯：mcp-server/src 導讀}
作為 AI 時代的新產物，MCP Server 的加入讓這個專案與傳統沙盒有了本質的區別：
\begin{itemize}
    \item \textbf{\texttt{index.ts}}：實作了 MCP 協定的 \texttt{ListToolsRequestSchema\allowbreak} 與 \texttt{CallToolRequestSchema\allowbreak}。目前對外暴露了 \texttt{translate\-error} 與 \texttt{suggest\-repair} 兩個工具。
    \item \textbf{\texttt{translator.ts}}：內建了一個 Mapping 字典，將諸如 \texttt{EPERM} 或 \texttt{SIGSYS} 結合對應的 Syscall Number，轉譯為一段具有上下文的英文指導語（因為主流 LLM 在英文的理解上最為精確）。
\end{itemize}

\subsection{13.4 專案總結與學習反思}
透過這個專案，我們深入了解了 Linux Kernel 的底層運作機制。從 \texttt{chroot} 的不安全性，到 Namespaces 的強大；從 \texttt{ptrace} 的效能瓶頸，到 Seccomp User Notification 的優雅。
最重要的是，我們跨越了純粹的「系統程式設計」範疇，將底層的 C/Rust 系統與前端的 React/Node.js、甚至是最尖端的 AI MCP 協定結合在一起，創造出了一個真正具備實用價值的「智慧型基礎設施」。這不僅是一次技術上的挑戰，更是對未來軟體開發典範（AI Native Development）的一次深度實踐。


\end{document}
